% ==================== SECTION VALIDATION ====================

\element{validation}{
    \section{Fonction de calcul sécurisée avec gestion d'erreurs}

    En électronique, le calcul de la résistance d'un diviseur de tension nécessite de valider les données pour éviter des erreurs critiques.

    \textbf{Rappel :} Dans un diviseur de tension, si $U_{total} = 12V$ et on veut $U_{sortie} = 5V$ avec $R_2 = 1000\Omega$, alors :
    $$R_1 = R_2 \times \frac{U_{total} - U_{sortie}}{U_{sortie}}$$

    \begin{question}{valid-1}
        Écrire une fonction complète :

        \lstinline[language=C]|float calculer_R1(float u_total, float u_sortie, float r2)|

        Cette fonction doit :
        \begin{enumerate}
            \item \textbf{Vérifier} que toutes les valeurs sont strictement positives
            \item \textbf{Vérifier} que $U_{sortie} < U_{total}$ (sinon le diviseur ne fonctionne pas)
            \item \textbf{Vérifier} que $U_{sortie} \neq 0$ (division par zéro)
            \item Si \textbf{une erreur est detectee} : retourner \lstinline[language=C]|-1.0|
            \item Si \textbf{tout est correct} : calculer et retourner $R_1$
        \end{enumerate}

        \textbf{Écrire uniquement la fonction complète avec tous les tests de validation (pas besoin du main).}

        \evaluationProfGrille{2}{13}
    \end{question}

    \begin{question}{valid-2}
        \textbf{Question bonus}\\
        Ecrire un programme \lstinline[language=C]|main| qui teste la fonction avec au moins \textbf{3 cas differents} :
        \begin{itemize}
            \item 1 cas valide qui doit réussir
            \item 2 cas invalides qui doivent échouer
        \end{itemize}

        Pour chaque test, afficher si le calcul a réussi ou échoué.

        \textbf{Exemple d'utilisation attendue :}

        \lstinputlisting[language=C]{sources/codes/exemple_validation.c}

        \textbf{Écrire le programme main complet (avec les \#include nécessaires).}

        \evaluationProfGrille{2}{16}
    \end{question}
}
